% GENERATED FILE. Edit canonical lessons, then run npm run generate:books.
% canonical-source-hash: fnv1a64:3155a7a75f6b86f1
% canonical-lessons: PT-C485-arranhar, PT-C485-arranhao, PT-C485-nodoa-negra, PT-C485-inchar, PT-C485-inchado

\chapter{Arranhões e nódoas negras}
\label{ch:pt-arranhoes-e-nodoas-negras}
\hlchaptermodality{485}

\begin{chapteropening}
\textbf{By the end of this chapter:} \emph{I can describe small injuries.}

Complete the last lesson of chapter 485: I can describe small injuries.
\end{chapteropening}

\section[arranhar]{arranhar — to scratch (and leave a mark)}

\label{lesson:PT-C485-arranhar}



\begin{quote}
Before the new one: say the Portuguese for nausea, then the Portuguese for to vomit.
\end{quote}

\subsection*{You'll want to know: arranhar}
\textbf{arranhar} — \textquotedblleft{}to scratch (and leave a mark)\textquotedblright{}.

\begin{quote}
O gato arranhou-me a mão.
\end{quote}

\begin{quote}
\emph{The cat scratched my hand.}
\end{quote}

\emph{Coçar} is for an itch; \emph{arranhar} leaves a mark.

\begin{cousinweb}
\emph{arranhar} is an old Portuguese verb (EtymDB), but its origin is not settled.
\end{cousinweb}

\begin{grammarlens}[title={What you've built}]
One more word for telling, explaining and describing at B1.
\end{grammarlens}

\subsection*{Your turn}
\begin{itemize}
\raggedright
  \item \emph{Say it:} \emph{arranhar}
  \item \emph{Say it:} \emph{arranhar}, once more
  \item \emph{Say it:} a sentence of your own with \emph{arranhar}
  \item \emph{Recall:} say the Portuguese for nausea, then the Portuguese for to vomit, then say \emph{arranhar} again
\end{itemize}

\begin{tcolorbox}[breakable,colback=teal!4,colframe=teal!35!black,title={Before you move on}]
What does arranhar mean? (\textquotedblleft{}To scratch (and leave a mark)\textquotedblright{}.) Say it once more.
\end{tcolorbox}
\section[o arranhão]{o arranhão — scratch (the mark)}

\label{lesson:PT-C485-arranhao}



\begin{quote}
Before the new one: say the Portuguese for to vomit, then the Portuguese for to scratch.
\end{quote}

\subsection*{You'll want to know: o arranhão}
\textbf{o arranhão} — \textquotedblleft{}scratch (the mark)\textquotedblright{}.

\begin{quote}
Tenho um arranhão no braço.
\end{quote}

\begin{quote}
\emph{I have a scratch on my arm.}
\end{quote}

Built on \emph{arranhar}.

\begin{cousinweb}
\emph{arranhão} is built on \emph{arranhar} \textquotedblleft{}to scratch\textquotedblright{} with \emph{-ão}; the origin of the verb is not settled.
\end{cousinweb}

\begin{grammarlens}[title={What you've built}]
One more word for telling, explaining and describing at B1.
\end{grammarlens}

\subsection*{Your turn}
\begin{itemize}
\raggedright
  \item \emph{Say it:} \emph{o arranhão}
  \item \emph{Say it:} \emph{o arranhão}, once more
  \item \emph{Say it:} a sentence of your own with \emph{o arranhão}
  \item \emph{Recall:} say the Portuguese for to vomit, then the Portuguese for to scratch, then say \emph{o arranhão} again
\end{itemize}

\begin{tcolorbox}[breakable,colback=teal!4,colframe=teal!35!black,title={Before you move on}]
What does o arranhão mean? (\textquotedblleft{}Scratch (the mark)\textquotedblright{}.) Say it once more.
\end{tcolorbox}
\section[a nódoa negra]{a nódoa negra — bruise}

\label{lesson:PT-C485-nodoa-negra}



\begin{quote}
Before the new one: say the Portuguese for to scratch, then the Portuguese for scratch.
\end{quote}

\subsection*{You'll want to know: a nódoa negra}
\textbf{a nódoa negra} — \textquotedblleft{}bruise\textquotedblright{}.

\begin{quote}
Bati com a perna na cama e fiquei com uma nódoa negra.
\end{quote}

\begin{quote}
\emph{I hit my leg on the bed and got a bruise.}
\end{quote}

Literally \textquotedblleft{}a black stain\textquotedblright{}. In Brazil: \emph{a mancha roxa}.

\begin{cousinweb}
\emph{nódoa} is inherited from Latin \emph{notula} \textquotedblleft{}a little mark\textquotedblright{} (EtymDB), a diminutive of \emph{nota} \textquotedblleft{}a mark\textquotedblright{}; \emph{negra} comes from Latin \emph{niger} \textquotedblleft{}black\textquotedblright{} (Lewis \& Short).
\end{cousinweb}

\begin{grammarlens}[title={What you've built}]
One more word for telling, explaining and describing at B1.
\end{grammarlens}

\subsection*{Your turn}
\begin{itemize}
\raggedright
  \item \emph{Say it:} \emph{a nódoa negra}
  \item \emph{Say it:} \emph{a nódoa negra}, once more
  \item \emph{Say it:} a sentence of your own with \emph{a nódoa negra}
  \item \emph{Recall:} say the Portuguese for to scratch, then the Portuguese for scratch, then say \emph{a nódoa negra} again
\end{itemize}

\begin{tcolorbox}[breakable,colback=teal!4,colframe=teal!35!black,title={Before you move on}]
What does a nódoa negra mean? (\textquotedblleft{}Bruise\textquotedblright{}.) Say it once more.
\end{tcolorbox}
\section[inchar]{inchar — to swell}

\label{lesson:PT-C485-inchar}



\begin{quote}
Before the new one: say the Portuguese for scratch, then the Portuguese for bruise.
\end{quote}

\subsection*{You'll want to know: inchar}
\textbf{inchar} — \textquotedblleft{}to swell\textquotedblright{}.

\begin{quote}
O tornozelo inchou depois da queda.
\end{quote}

\begin{quote}
\emph{The ankle swelled after the fall.}
\end{quote}

Regular: \emph{incha}, \emph{inchou}.

\begin{cousinweb}
\emph{inchar} is inherited, through Old Portuguese \emph{enchar}, from Latin \emph{inflāre} \textquotedblleft{}to blow into, to puff up\textquotedblright{} (EtymDB).
\end{cousinweb}

\begin{grammarlens}[title={What you've built}]
One more word for telling, explaining and describing at B1.
\end{grammarlens}

\subsection*{Your turn}
\begin{itemize}
\raggedright
  \item \emph{Say it:} \emph{inchar}
  \item \emph{Say it:} \emph{inchar}, once more
  \item \emph{Say it:} a sentence of your own with \emph{inchar}
  \item \emph{Recall:} say the Portuguese for scratch, then the Portuguese for bruise, then say \emph{inchar} again
\end{itemize}

\begin{tcolorbox}[breakable,colback=teal!4,colframe=teal!35!black,title={Before you move on}]
What does inchar mean? (\textquotedblleft{}To swell\textquotedblright{}.) Say it once more.
\end{tcolorbox}
\section[inchado]{inchado — swollen}

\label{lesson:PT-C485-inchado}



\begin{quote}
Before the new one: say the Portuguese for bruise, then the Portuguese for to swell.
\end{quote}

\subsection*{You'll want to know: inchado}
\textbf{inchado} — \textquotedblleft{}swollen\textquotedblright{}.

\begin{quote}
Tenho o joelho inchado.
\end{quote}

\begin{quote}
\emph{My knee is swollen.}
\end{quote}

The participle of \emph{inchar}.

\begin{cousinweb}
\emph{inchado} is the participle of \emph{inchar} \textquotedblleft{}to swell\textquotedblright{}, inherited from Latin \emph{inflāre} \textquotedblleft{}to puff up\textquotedblright{} (EtymDB).
\end{cousinweb}

\begin{grammarlens}[title={What you've built}]
One more word for telling, explaining and describing at B1.
\end{grammarlens}

\subsection*{Your turn}
\begin{itemize}
\raggedright
  \item \emph{Say it:} \emph{inchado}
  \item \emph{Say it:} \emph{inchado}, once more
  \item \emph{Say it:} a sentence of your own with \emph{inchado}
  \item \emph{Recall:} say the Portuguese for bruise, then the Portuguese for to swell, then say \emph{inchado} again
\end{itemize}

\begin{tcolorbox}[breakable,colback=teal!4,colframe=teal!35!black,title={Before you move on}]
What does inchado mean? (\textquotedblleft{}Swollen\textquotedblright{}.) Say it once more.
\end{tcolorbox}
